\chapter{Training Experiments}


We performed the training of our policies in our own NVIDIA GeForce GPUs available in our own lab, we then let the training run while randomizing the starting states (positions and velocities). Initial policies were trained with relatively myopic policies (gamma = 0.991) and a capture radius of 0.75 meters. The hyperparameters utilized throughout these experimetns  are described in appendix \ref{appendix:hyperparameters}. Do note that tested on a number of different hyperparameters before converging on these, so their choice for the trainings in this work is empirical.


We chose to focus on three different test cases for this thesis, where we just varied the actuation limit $u_{\max}$, the max velocity $v_{\max}$, and the max initial velocity in any given episode $v_{\max}^{\mathrm{IC}}$, leaving all other training hyperparameters equivalent between the trhee runs. We run these with a KF off for a full observability run, and with a KF on for a partially observable run via behavior cloning. The KF on results have a (NUMBER needed) noise on the prior adversary action.

The three test cases that we'll focus on in the remainder of this work are:


\begin{itemize}
    \item Training case 1: $u_{\max} = 0.1~\mathrm{m/s^2}$, $v_{\max} = 1.0~\mathrm{m/s}$, $v_{\max}^{\mathrm{IC}} = 0.05~\mathrm{m/s}$
    \item Training case 2: $u_{\max} = 0.5~\mathrm{m/s^2}$, $v_{\max} = 1.5~\mathrm{m/s}$, $v_{\max}^{\mathrm{IC}} = 0.25~\mathrm{m/s}$
    \item Training case 3: $u_{\max} = 2.0~\mathrm{m/s^2}$, $v_{\max} = 1.0~\mathrm{m/s}$, $v_{\max}^{\mathrm{IC}} = 1.0~\mathrm{m/s}$
\end{itemize}

These values are meant to test the PARL+BC-produced policies and their corresponding Monte Carlo evaluations at different actuation limits and max velocity saturation limits.


\section{Trainings Curves}
\label{sec:Net_Training_Curves}

\section{Training case 1: $u_{\max} = 0.1~\mathrm{m/s^2}$, $v_{\max} = 1.0~\mathrm{m/s}$, $v_{\max}^{\mathrm{IC}} = 0.05~\mathrm{m/s}$}

\Needspace{0.45\textheight}
\begin{figure}[H]
\centering
\begin{minipage}[t]{0.48\textwidth}
\centering
\textbf{KF off}\\[0.25em]
\includegraphics[width=\linewidth]{\detokenize{Figures/Merged_Thesis/Training_Policy_0.1u_1vmax_0.05_icVmax/combined_training_curve_kf_off.png}}
\end{minipage}
\hfill
\begin{minipage}[t]{0.48\textwidth}
\centering
\textbf{KF on}\\[0.25em]
\includegraphics[width=\linewidth]{\detokenize{Figures/Merged_Thesis/Training_Policy_0.1u_1vmax_0.05_icVmax/combined_training_curve_kf_on.png}}
\end{minipage}
\caption{Combined training return curves for $u_{\max} = 0.1~\mathrm{m/s^2}$, $v_{\max} = 1.0~\mathrm{m/s}$, $v_{\max}^{\mathrm{IC}} = 0.05~\mathrm{m/s}$, across 5 interchanging training phases.}
\label{fig:mc-training-policy-0-1u-1vmax-0-05-icvmax-training-curves-kf}
\end{figure}


\section{Training case 2: $u_{\max} = 0.5~\mathrm{m/s^2}$, $v_{\max} = 1.5~\mathrm{m/s}$, $v_{\max}^{\mathrm{IC}} = 0.25~\mathrm{m/s}$}

\Needspace{0.45\textheight}
\begin{figure}[H]
\centering
\begin{minipage}[t]{0.48\textwidth}
\centering
\textbf{KF off}\\[0.25em]
\includegraphics[width=\linewidth]{\detokenize{Figures/Merged_Thesis/Training_Policy_0.5u_1.5vmax_0.25_icVmax/combined_training_curve_kf_off.png}}
\end{minipage}
\hfill
\begin{minipage}[t]{0.48\textwidth}
\centering
\textbf{KF on}\\[0.25em]
\includegraphics[width=\linewidth]{\detokenize{Figures/Merged_Thesis/Training_Policy_0.5u_1.5vmax_0.25_icVmax/combined_training_curve_kf_on.png}}
\end{minipage}
\caption{Combined training return curves for $u_{\max} = 0.5~\mathrm{m/s^2}$, $v_{\max} = 1.5~\mathrm{m/s}$, $v_{\max}^{\mathrm{IC}} = 0.25~\mathrm{m/s}$, across 5 interchanging training phases.}
\label{fig:mc-training-policy-0-5u-1-5vmax-1-0-icvmax-training-curves-kf}
\end{figure}


\section{Training Case 3: $u_{\max} = 2.0~\mathrm{m/s^2}$, $v_{\max} = 1.0~\mathrm{m/s}$, $v_{\max}^{\mathrm{IC}} = 1.0~\mathrm{m/s}$}

\Needspace{0.45\textheight}
\begin{figure}[H]
\centering
\begin{minipage}[t]{0.48\textwidth}
\centering
\textbf{KF off}\\[0.25em]
\includegraphics[width=\linewidth]{\detokenize{Figures/Merged_Thesis/Training_Policy_2.0u_1vmax_1.0_icVmax/combined_training_curve_kf_off.png}}
\end{minipage}
\hfill
\begin{minipage}[t]{0.48\textwidth}
\centering
\textbf{KF on}\\[0.25em]
\includegraphics[width=\linewidth]{\detokenize{Figures/Merged_Thesis/Training_Policy_2.0u_1vmax_1.0_icVmax/combined_training_curve_kf_on.png}}
\end{minipage}
\caption{Combined training return curves for $u_{\max} = 2.0~\mathrm{m/s^2}$, $v_{\max} = 1.0~\mathrm{m/s}$, $v_{\max}^{\mathrm{IC}} = 1.0~\mathrm{m/s}$, across 5 interchanging training phases.}
\label{fig:mc-training-policy-2-0u-1vmax-1-0-icvmax-training-curves-kf}
\end{figure}

\subsection{Training Analysis}

Overall, just from analyzing each curve described in figures \ref{fig:mc-training-policy-0-1u-1vmax-0-05-icvmax-training-curves-kf}, \ref{fig:mc-training-policy-0-5u-1-5vmax-1-0-icvmax-training-curves-kf}, and \ref{fig:mc-training-policy-2-0u-1vmax-1-0-icvmax-training-curves-kf} we see that PARL+BC framework manages to produce strong, convergent policies across the multiple phases of training. Furthermore, the partially observable results follow loosely the same pattern as the fully observable ones.

By studying the mean reward at which the polciies converge, we see that the defender manages to get slightly stronger policies than that of the attacker as the phases progress and, consequently, the attacker policies become slightly weaker. This is a good sign, as it signals that the defender is progressively stronger and more generalizable thanks to the opponent randomization strategy that was described before, allowing it to be resilient throughout the progression of the training.

 

